\section{模型评价与推广}

\subsection{模型优点}
\begin{enumerate}[leftmargin=2em]
  \item \textbf{口径统一、数据一致继承}：四问共用附录 2 时间/能耗规则与附录 3 通信规则，对象编号、单位与计算口径完全一致，问题一的载荷/组批、问题二的调度、问题三的通信、问题四的分区逐层继承，模型体系自洽。
  \item \textbf{多方案与层级求解}：问题一比较线性加权 MILP、Chebyshev MILP 和 NSGA-II/Minkowski 三种组批方法；问题二比较全局 NSGA-II 与贪心基础方案；问题三、四分别采用高复杂度主解与基础解，算法角色和比较口径清晰。
  \item \textbf{结果可追溯且分项核验}：各主链方案均通过货箱全覆盖（$80/80$）；问题三两种方案通信缺口均为 $0$。问题二全局 NSGA-II 硬约束违反数为 $0$，但问题二基础方案有 $5$ 个违反，问题三主解和基础解分别有 $12$、$18$ 个违反，论文对此不再笼统宣称“所有方案完全可行”。
  \item \textbf{权衡关系清晰}：每问均以多目标形式给出并量化权衡（架次—能耗—时间、及时性—完工—能耗—架次、运输—中继代价、组数—规模—均衡—缺口），为应急决策提供可选前沿而非单一答案。
\end{enumerate}

\subsection{模型缺点}
\begin{enumerate}[leftmargin=2em]
  \item 航段采用两节点水平直线近似，未刻画绕障与风场影响，复杂地形下能耗可能被低估；
  \item 通信状态按离散相位采样判定，采样粒度与真实连续视线之间存在细微误差；
  \item 问题二、三的高维调度依赖启发式，NSGA-II 代表解为近似 Pareto 最优，不保证全局最优。
\end{enumerate}

\subsection{改进方向与推广}

针对上述缺点，可引入三维路径规划（A*/RRT）替代直线航段以刻画绕障，用连续时间视线积分细化通信判定，并以列生成或分支定价提升大规模调度的精确性。本文的“组批—多行程调度—通信协同—任务分区”建模框架可直接推广至山火、地震等其他灾种的应急无人机保障，以及城市低空物流、海岛补给、电力巡检等需要运输与通信协同、资源受限调度的场景。
